\chapter{SAN design}
\label{ch:san-design}

\section{Purpose}
\label{purpose}

There are a few primary patterns of storage provided by SAN:

\begin{description}
\item[Online host-specific block storage]
essentially replacing directly-attached block devices on the hosts.
\item[Online clustered block storage]
allowing a number of hosts to share a block device that is formatted
with a clustered/shared file system.
\item[Offline block storage]
replacing removeable drives on hosts. This may be for archival purposes.
\item[Backup]
where the SAN facilitates access to block devices for backup. These may
be disk-backed LUNs on storage appliance. May also include tape-based
backup systems on SAN-attached tape equipment.
\end{description}

Some additional host-specific questions:

\begin{itemize}
\item
  Will hosts be booting from the SAN? (There are a number of ways to
  provision this.)
\item
  Are the hosts
\end{itemize}

\section{SAN type}\label{san-type}

A major decision is whether to deploy FC or IP SAN.

\begin{description}
\item[Motivation for IP SAN]
Low/no-cost entry; leverage knowledge; re-use of infrastructure;
multi-site capable if IP routed.
\item[Motivation for FC SAN]
Performance; reliability; inherent SAN/LAN demarcation.
\end{description}

Other notes:

\begin{itemize}
\item
  Problem-solving technologies available: FCoE, FCoIP, iSCSI-FC
  bridging. Unwise to base network around these technologies.
\item
  Depending on organisation structure and culture, there may be
  different teams responsible for storage vs general-purpose networking.
\end{itemize}

\section{High Availability}\label{high-availability}

SANs hold high-value information critical to continuity of service. A
naively designed SAN may exhibit worse availability characteristics than
simple direct-attached storage on individual hosts.

\subsection{Storage layer
availability}\label{storage-layer-availability}

Storage appliances normally include a number of features to ensure high
availability, .

\begin{description}
\item[Disk redundancy]
where volumes / LUNs are provisioned such that the array is redundant
against disk failures. May include RAID, hot spares etc.
\item[Dual power-supply]
where each storage applicance has two Power Supply Units (PSU). Each PSU
independently capable of supporting the load.

\begin{itemize}

\item
  Data centre environment should have dual power supply to take full
  advantage of dual PSU.
\end{itemize}
\item[Redundant cooling components]
permitting the storage appliance to continue operating despite the
failure of 1 (or sometimes more) fans. Should be easily replaced. Often
hot swappable.
\item[Dual controller]
where the storage array has two controller modules
(i.e.~high-performance embedded computer). This normally works so that
either controller can independently provide all services. May be hot
swappable, .
\item[Remote management capabilities]
that allow the storage appliance to alert administrators should any
failures occur. Essential that all alerts are enabled and acted upon to
avoid storage appliance remaining in a degraded state (e.g.~failed PSU)
as redundancy will hide this situation.
\end{description}

\subsection{Fabric layer availability}\label{fabric-layer-availability}

The fabric (regardless of type) can represent a significant single point
of failure. Consider the simple SAN layout in .

This is vulnerable to a failure of the switch SW1. Duplicating the SAN
fabric averts this, .

SAN fabric components (e.g.~switches) should be dual-powered where
possible. This may be via dual-PSU or via external transfer switches.

\subsection{Host layer availability}\label{host-layer-availability}

Assuming that there is redundancy in the fabric, the host layer requires
approriate redundancy to connect to the fabric.

To transparently utilise the redundant fabric we must ensure that
multipathing is correctly configured.

\subsection{Data centre environment}\label{data-centre-environment}

A SAN will normally be provisioned within a data centre environment. We
should assume at a minimum:

\begin{description}
\item[Power supply]
is protected by UPS, possibly generator and has diverse A and B power
paths.
\item[Cooling]
is sufficient to meet the load and incorporates redundancy to maintain
operation in the event of a failure.
\item[Fire detection and suppression]
to protect data and ensure service availability.
\item[Physical security]
appropriate to the data and service risk profile.
\end{description}

The provision of a SAN should dictate that the infrastructural services
are sufficient.
